\begin{ouverture}[vecteursplan]{10}{Droites du plan}{Partie 3 · Géométrie}
{Équations · Pentes · Parallélisme · Intersections}
\objectif{Passer d’une droite à ses équations et réciproquement.}
\objectif{Construire une droite à partir de données variées.}
\objectif{Résoudre une intersection et interpréter un système.}
\prerequis{Vecteurs, déterminants et fonctions affines.}
\end{ouverture}
\pagestyle{otMain}\essentiel\fontsize{11}{15}\selectfont
\essSection{1}{Deux écritures pour une droite}
\begin{definition}Une équation cartésienne de droite est $ax+by+c=0$, avec $(a;b)\ne(0;0)$. Un vecteur directeur non nul est $(-b;a)$. Il décrit une direction parallèle à la droite.\end{definition}
\begin{propriete}Si $b\ne0$, l’équation devient $y=mx+p$, avec $m=-a/b$ et $p=-c/b$. Si $b=0$, la droite est verticale : $x=-c/a$. Elle n’a pas de pente et ne représente pas une fonction de $x$.\end{propriete}
\begin{savoirfaire}[À partir de deux points]
Pour $A$ et $B$ distincts avec $x_A\ne x_B$, calculer $m=(y_B-y_A)/(x_B-x_A)$ puis $p=y_A-mx_A$. Si leurs abscisses sont égales, écrire directement $x=x_A$.
\end{savoirfaire}
\begin{savoirfaire}[À partir d’un directeur]
Si $d$ passe par $A$ et a pour directeur $\vec u=(u;v)\ne\vec0$, un point $M(x;y)$ est sur $d$ si $\det(\vec u,\vect{AM})=0$, soit $u(y-y_A)-v(x-x_A)=0$.
\end{savoirfaire}
\clearpage
\essSection{2}{Positions relatives et systèmes}
\begin{propriete}Deux droites non verticales de même pente sont parallèles : confondues si leurs ordonnées à l’origine sont aussi égales, distinctes sinon. Avec des directeurs, le déterminant nul caractérise le parallélisme.\end{propriete}
\begin{savoirfaire}[Intersection]
Pour $y=m_1x+p_1$ et $y=m_2x+p_2$, résoudre $(m_1-m_2)x=p_2-p_1$, puis calculer $y$ dans l’une des équations. Vérifier le point dans les deux équations initiales.
\end{savoirfaire}
\begin{exemple}$y=2x+1$ et $y=-x+7$ se coupent pour $3x=6$, donc au point $(2;5)$. Deux pentes distinctes donnent exactement une intersection.\end{exemple}
\begin{attention}Une équation peut être multipliée par un réel non nul sans changer la droite. Les écritures $x-2y+3=0$ et $2x-4y+6=0$ représentent ainsi la même droite.\end{attention}
\otIntertitre{Avec un logiciel de géométrie}
Créer deux points, tracer leur droite et afficher son équation. Déplacer un point jusqu’à obtenir une droite verticale : observer ce qui change dans l’équation réduite. Expliquer ensuite ce changement avec les coordonnées.
\clearpage\rubrique{Découverte}
\begin{blocLivre}{ex-10-1}
\begin{exercice}[1]{Tester des points}
La droite $d$ a pour équation $2x-y+3=0$. Les points $A(1;5)$ et $B(2;6)$ lui appartiennent-ils ?
\label{fig-ex-10-1}
\figureGeo{droite-01}
\end{exercice}
\end{blocLivre}
\begin{blocLivre}{ex-10-2}
\begin{exercice}[1]{De l’équation au tracé}
Écrire $3x-2y+6=0$ sous forme réduite. Donner deux points de la droite.
\label{fig-ex-10-2}
\figureGeo{droite-02}
\end{exercice}
\end{blocLivre}
\begin{blocLivre}{ex-10-3}
\begin{exercice}[1]{Une droite verticale}
Quelle droite représente $x=2$ ? Donner un vecteur directeur non nul et expliquer pourquoi on ne peut pas l’écrire $y=mx+p$.
\label{fig-ex-10-3}
\figureGeo{droite-03}
\end{exercice}
\end{blocLivre}
\begin{blocLivre}{ex-10-4}
\begin{exercice}[1]{Lire une pente}
Une droite passe par $A(0;-1)$ et $B(3;5)$. Calculer sa pente et son équation réduite.
\label{fig-ex-10-4}
\figureGeo{droite-04}
\end{exercice}
\end{blocLivre}
\begin{blocLivre}{ex-10-5}
\begin{exercice}[1]{Un point et une pente}
Donner une équation de la droite de pente $-2$ passant par $A(1;3)$.
\label{fig-ex-10-5}
\figureGeo{droite-05}
\end{exercice}
\end{blocLivre}

\clearpage\rubrique{Entraînement}
\begin{blocLivre}{ex-10-6}
\begin{exercice}[2]{Un point et un directeur}
Déterminer une équation de la droite passant par $A(2;1)$ et de directeur $(3;-2)$.
\label{fig-ex-10-6}
\figureGeo{droite-06}
\end{exercice}
\end{blocLivre}
\begin{blocLivre}{ex-10-7}
\begin{exercice}[2]{Deux points quelconques}
Déterminer l’équation réduite de $(AB)$ pour $A(-1;4)$ et $B(3;-2)$.
\label{fig-ex-10-7}
\figureGeo{droite-07}
\end{exercice}
\end{blocLivre}
\begin{blocLivre}{ex-10-8}
\begin{exercice}[2]{Parallèles ou confondues ?}
Comparer les droites $d_1:y=2x-1$, $d_2:4x-2y-2=0$ et $d_3:y=2x+3$.
\label{fig-ex-10-8}
\figureGeo{droite-08}
\end{exercice}
\end{blocLivre}
\begin{blocLivre}{ex-10-9}
\begin{exercice}[2]{Une intersection}
Déterminer l’intersection de $d:y=2x+1$ et $e:y=-x+7$.
\label{fig-ex-10-9}
\figureGeo{droite-09}
\end{exercice}
\end{blocLivre}
\begin{blocLivre}{ex-10-10}
\begin{exercice}[2]{Intersection avec une verticale}
Trouver le point commun à $x=-2$ et $3x+2y=8$.
\label{fig-ex-10-10}
\figureGeo{droite-10}
\end{exercice}
\end{blocLivre}

\clearpage\rubrique{Problèmes}
\begin{blocLivre}{ex-10-11}
\begin{exercice}[3]{Des tarifs et des droites}
Deux livreurs facturent respectivement $A(d)=5+2d$ et $B(d)=11+d$, pour une distance $d\ge0$ en kilomètres. Comparer leurs tarifs selon la distance.
\label{fig-ex-10-11}
\figureGeo{droite-11}
\end{exercice}
\end{blocLivre}
\begin{blocLivre}{ex-10-12}
\begin{exercice}[3]{Une famille de droites}
Pour $m\in\mathbb R$, on considère $d_m:y=mx+2-m$. Montrer qu’elles passent toutes par un même point. Laquelle est parallèle à $y=3x-5$ ?
\label{fig-ex-10-12}
\figureGeo{droite-12}
\end{exercice}
\end{blocLivre}
\begin{blocLivre}{ex-10-13}
\begin{exercice}[3]{Un alignement à deux méthodes}
$A(-2;1)$, $B(2;3)$ et $C(6;5)$. Montrer leur alignement par une équation de droite, puis par un déterminant.
\label{fig-ex-10-13}
\figureGeo{droite-13}
\end{exercice}
\end{blocLivre}
\begin{blocLivre}{ex-10-14}
\begin{exercice}[3]{Un système à interpréter}
Discuter les solutions des systèmes formés de $x+2y=4$ avec, d’une part, $2x+4y=8$, et d’autre part $2x+4y=10$.
\label{fig-ex-10-14}
\figureGeo{droite-14}
\end{exercice}
\end{blocLivre}
\begin{blocLivre}{ex-10-15}
\begin{exercice}[3]{Une médiane à déterminer}
$A(0;4)$, $B(-2;0)$ et $C(4;0)$. Déterminer une équation de la médiane issue de $A$, puis son intersection avec la droite $y=2$.
\label{fig-ex-10-15}
\figureGeo{droite-15}
\end{exercice}
\end{blocLivre}

\clearpage
